\subsection{社区服务站办理质量对照}
本例是一份虚构的社区公共服务月报，对照东城、南园、西河、北湖四个服务站在十二类事项上的办理质量。
按期率是按期办结件数占本期办结件数的比例；回访率是完成回访件数占本期办结件数的比例，两个指标独立展示。
第一层表头标出服务站，第二层区分按期率与回访率；所有指标共用 $[0,1]$ 范围，显示三位小数。
例如 0.960 表示 96.0\%；空白表示本期尚未统计该项，0.000 表示该项有办结记录，但没有按期办结或完成回访的记录，两者不能互换。
本例只演示排版，不用于评价真实服务站。

中文独立文档使用 ctexart 与 XeLaTeX；导言区加载 \pkg{tabularray} 和 \pkg{gradbars}：
\begin{codeonly}
\usepackage{gradbars}
\usepackage{tabularray}
\end{codeonly}
\key{\gradbar} 绘制固定 8\,mm 的短条，三位小数显示在条形右侧。
声明中的 \key{empty=skip} 跳过空白；图形处理仅作用于第三行以后的指标区域，两个文字表头不会传入数值入口。

\begin{example}[text and listing,sidebyside=false,breakable]{社区服务月报：四站对照与两级表头}
\begingroup
\gradbarsnewcommand[type=bar,args={m},data={#1},empty=skip,
  theme=gray,min=0,max=1,width=8mm,precision=3,baseline=text,
  label=outside,label width=6mm,label gap=1mm]{\metricbar}
\begin{center}
\begin{tblr}{
  colspec={Q[l,m]*{8}{Q[c,m]}},
  colsep=0.8mm,rowsep=1pt,
  cells={font=\footnotesize},
  cell{1}{even}={c=2}{c},
  cell{3-Z}{2-Z}={cmd=\metricbar},
  hline{1,Z}={0.08em},hline{3}={0.05em},
  hline{2}={2-3}{0.03em,endpos,lr=-1},
  hline{2}={4-5}{0.03em,endpos,lr=-1},
  hline{2}={6-7}{0.03em,endpos,lr=-1},
  hline{2}={8-9}{0.03em,endpos,lr=-1}
}
         & 东城站 & & 南园站 & & 西河站 & & 北湖站 & \\
服务事项 & 按期率 & 回访率 & 按期率 & 回访率
     & 按期率 & 回访率 & 按期率 & 回访率 \\
证件办理 & 0.960 & 0.720 & 0.880 & 0.450
         & 0.910 & 0.630 & 0.940 & 0.810 \\
户籍咨询 & 0.840 & 0.380 & 0.930 & 0.760
         & 0.790 & 0.520 & 0.900 & 0.640 \\
养老服务 & 0.980 & 0.920 & 0.870 & 0.610
         &  &  & 0.950 & 0.840 \\
医疗咨询 & 0.760 & 0.490 & 0.890 & 0.680
         & 0.920 & 0.800 &  &  \\
就业登记 & 0.830 & 0.560 & 0.710 & 0.320
         &  &  & 0.880 & 0.740 \\
社保查询 & 0.910 & 0.670 & 0.950 & 0.830
         & 0.860 & 0.410 & 0.780 & 0.590 \\
住房咨询 & 1.000 & 0.940 & 0.900 & 0.730
         & 0.820 & 0.580 & 0.960 & 0.860 \\
档案查阅 & 0.730 & 0.240 & 0.850 & 0.470
         & 0.680 & 0.000 & 0.890 & 0.620 \\
便民预约 & 0.990 & 0.870 & 0.920 & 0.690
         & 0.740 & 0.280 & 0.860 & 0.510 \\
困难帮扶 & 0.880 & 0.780 & 0.960 & 0.910
         &  &  & 0.810 & 0.540 \\
法律咨询 & 0.690 & 0.310 & 0.820 & 0.570
         & 0.930 & 0.850 &  &  \\
志愿服务 & 0.900 & 0.650 & 0.770 & 0.430
         &  &  & 0.850 & 0.180 \\
\end{tblr}
\end{center}
\endgroup
\end{example}
\key{cell{1}{even}={c=2}{c}} 将第一行每个服务站标题横跨两列；
\key{cell{3-Z}{2-Z}} 只为数据区域指定绘图命令。
各列按内容自然分配宽度，列间留白为 0.8\,mm；\key{label=outside} 将数字放在条外，
\key{label width=6mm,label gap=1mm} 为数字和条形间隙预留位置。
表头保持常规字重，四组标题下各用一段细横线区分。
灰色数据条保留数字，便于黑白打印；条长表示原始比例，不额外转换为百分数。

\noindent 完整横向 A4 文档见 \texttt{docs/examples/}\allowbreak\texttt{precision-recall-tblr.tex}，从仓库根目录编译。
